% ============================================================
\chapter{Lois des Grands Nombres}
\label{ch:grands_nombres}
% ============================================================

Pourquoi la fr\'equence observ\'ee d'un \'ev\'enement se stabilise-t-elle autour de sa probabilit\'e th\'eorique quand on r\'ep\`ete l'exp\'erience un grand nombre de fois~? Cette question, au c\oe ur de la pens\'ee probabiliste, a re\c{c}u sa premi\`ere r\'eponse de Jacob Bernoulli en 1713, dans son \emph{Ars Conjectandi} publi\'e \`a titre posthume. Bernoulli montre que la fr\'equence empirique converge vers la probabilit\'e --- mais sa d\'emonstration est longue et laborieuse. Tchebyshev, en 1867, simplifie consid\'erablement la preuve gr\^ace \`a son in\'egalit\'e. Kolmogorov, en 1933, couronne l'\'edifice avec la \emph{loi forte des grands nombres}, qui garantit la convergence presque s\^ure. Ce chapitre retrace cette progression, des in\'egalit\'es de concentration aux diff\'erents modes de convergence stochastique.

\begin{intuition}
La loi des grands nombres est l'un des r\'esultats les plus fondamentaux de la
th\'eorie des probabilit\'es. Elle justifie l'interpr\'etation fr\'equentiste de
la probabilit\'e~: la fr\'equence empirique d'un \'ev\'enement converge vers sa
probabilit\'e th\'eorique lorsque le nombre d'exp\'eriences augmente.
\end{intuition}

% ------------------------------------------------------------
\section{Modes de convergence}
% ------------------------------------------------------------

\begin{definition}[Convergence en probabilit\'e]
Une suite $(X_n)$ \textbf{converge en probabilit\'e} vers $X$, not\'e
$X_n \xrightarrow{\PP} X$, si pour tout $\varepsilon > 0$~:
\[
    \lim_{n \to \infty} \PP(\abs{X_n - X} > \varepsilon) = 0.
\]
\end{definition}

\begin{definition}[Convergence presque s\^ure]
$(X_n)$ \textbf{converge presque s\^urement} vers $X$, not\'e
$X_n \xrightarrow{\text{p.s.}} X$, si~:
\[
    \PP\!\left(\lim_{n \to \infty} X_n = X\right) = 1,
\]
ou de mani\`ere \'equivalente, $\PP\!\left(\bigcap_{m=1}^{\infty}
\bigcup_{n=m}^{\infty} \{\abs{X_n - X} > \varepsilon\}\right) = 0$
pour tout $\varepsilon > 0$.
\end{definition}

\begin{definition}[Convergence dans $L^p$]
$(X_n)$ \textbf{converge dans $L^p$} ($p \geq 1$) vers $X$, not\'e
$X_n \xrightarrow{L^p} X$, si~:
\[
    \lim_{n \to \infty} \E[\abs{X_n - X}^p] = 0.
\]
\end{definition}

\begin{theoreme}[Hi\'erarchie des convergences]
\[
    X_n \xrightarrow{L^p} X \implies X_n \xrightarrow{\PP} X,
    \qquad
    X_n \xrightarrow{\text{p.s.}} X \implies X_n \xrightarrow{\PP} X.
\]
Les r\'eciproques sont fausses en g\'en\'eral. En revanche~:
$X_n \xrightarrow{\PP} X$ implique l'existence d'une sous-suite
$X_{n_k} \xrightarrow{\text{p.s.}} X$.
\end{theoreme}

\begin{attention}[title=\textbf{Convergence p.s.\ et convergence $L^p$}]
La convergence p.s.\ n'implique pas la convergence $L^p$, ni r\'eciproquement.
Chacune a ses propres conditions.
\end{attention}

% ------------------------------------------------------------
\section{Loi faible des grands nombres}
% ------------------------------------------------------------

\begin{theoreme}[Loi faible des grands nombres (LFGN)]
\label{thm:lfgn}
Soit $(X_n)_{n \geq 1}$ une suite de variables al\'eatoires
\textbf{ind\'ependantes et identiquement distribu\'ees} (i.i.d.) avec
$\E[X_1] = \mu$ et $\Var(X_1) = \sigma^2 < \infty$. Posons
$\bar{X}_n = \frac{1}{n}\sum_{i=1}^n X_i$. Alors~:
\[
    \bar{X}_n \xrightarrow{\PP} \mu, \quad
    \text{i.e.} \quad \forall \varepsilon > 0, \quad
    \lim_{n \to \infty} \PP(\abs{\bar{X}_n - \mu} > \varepsilon) = 0.
\]
\end{theoreme}

\begin{proof}
C'est une application directe de l'in\'egalit\'e de Bienaym\'e--Tchebychev.
\begin{align*}
    \PP(\abs{\bar{X}_n - \mu} > \varepsilon)
    &\leq \frac{\Var(\bar{X}_n)}{\varepsilon^2}
    = \frac{\Var\!\left(\frac{1}{n}\sum X_i\right)}{\varepsilon^2}
    = \frac{1}{n^2} \cdot \frac{n\sigma^2}{\varepsilon^2}
    = \frac{\sigma^2}{n\varepsilon^2}
    \xrightarrow[n \to \infty]{} 0. \qedhere
\end{align*}
\end{proof}

\begin{remarque}
La LFGN reste vraie sous l'hypoth\`ese plus faible que les $X_i$ sont non corr\'el\'ees
(pas n\'ecessairement ind\'ependantes) avec les m\^emes deux premiers moments.
\end{remarque}

\begin{intuition}[title=\textbf{Signification de la LFGN}]
La moyenne empirique $\bar{X}_n$ se concentre autour de $\mu$ quand $n$ est grand.
C'est le fondement th\'eorique des m\'ethodes de Monte Carlo et de l'estimation
statistique.
\end{intuition}

% ------------------------------------------------------------
\section{Loi forte des grands nombres}
% ------------------------------------------------------------

\begin{theoreme}[Loi forte des grands nombres (LFGN forte)]
\label{thm:lf_forte}
Soit $(X_n)_{n \geq 1}$ une suite i.i.d.\ avec $\E[\abs{X_1}] < \infty$ et
$\E[X_1] = \mu$. Alors~:
\[
    \bar{X}_n \xrightarrow{\text{p.s.}} \mu, \quad \text{i.e.} \quad
    \PP\!\left(\lim_{n \to \infty} \bar{X}_n = \mu\right) = 1.
\]
\end{theoreme}

\begin{remarque}
La loi forte ne n\'ecessite que l'existence du moment d'ordre 1 (pas de variance
finie). Sa d\'emonstration compl\`ete utilise des outils plus avanc\'es (troncation,
lemme de Borel--Cantelli).
\end{remarque}

\begin{proof}[Preuve dans le cas $\E\lbrack X_1^4\rbrack < \infty$]
Nous donnons la preuve sous l'hypoth\`ese plus forte $\E[X_1^4] < \infty$.
Sans perte de g\'en\'eralit\'e, supposons $\mu = 0$. Posons $S_n = \sum_{i=1}^n X_i$.
\begin{align*}
    \E[S_n^4] &= \E\!\left[\left(\sum_{i=1}^n X_i\right)^4\right]
    = \sum_{i} \E[X_i^4] + \binom{4}{2}\sum_{i \neq j} \E[X_i^2]\E[X_j^2]
\end{align*}
(les termes crois\'es d'ordre impair s'annulent par ind\'ependance et $\E[X_i] = 0$).
Donc $\E[S_n^4] = n\E[X_1^4] + 3n(n-1)(\E[X_1^2])^2 \leq Cn^2$ pour une constante $C$.
Par l'in\'egalit\'e de Markov~:
\[
    \PP\!\left(\abs{\bar{X}_n} > \varepsilon\right)
    = \PP\!\left(\abs{S_n} > n\varepsilon\right)
    \leq \frac{\E[S_n^4]}{n^4\varepsilon^4}
    \leq \frac{C}{n^2\varepsilon^4}.
\]
Comme $\sum_n 1/n^2 < \infty$, le premier lemme de Borel--Cantelli donne
$\PP(\limsup\{\abs{\bar{X}_n} > \varepsilon\}) = 0$ pour tout $\varepsilon > 0$,
d'o\`u $\bar{X}_n \to 0$ p.s.
\end{proof}

% ------------------------------------------------------------
\section{Applications}
% ------------------------------------------------------------

\subsection{Interpr\'etation fr\'equentiste}

\begin{corollaire}
Soit $A$ un \'ev\'enement de probabilit\'e $p$. Si l'on r\'ep\`ete l'exp\'erience
$n$ fois ind\'ependamment, la fr\'equence empirique
$\hat{p}_n = \frac{1}{n}\sum_{i=1}^n \mathbf{1}_{A_i}$ converge presque s\^urement
vers $p$.
\end{corollaire}

\subsection{M\'ethode de Monte Carlo}

\begin{exemple}[Estimation de $\pi$]
Soit $(X_i, Y_i)_{i \geq 1}$ i.i.d.\ uniformes sur $[0,1]^2$. Posons
$Z_i = \mathbf{1}_{X_i^2 + Y_i^2 \leq 1}$. Alors $\E[Z_i] = \pi/4$ et par la
loi des grands nombres~:
\[
    \frac{4}{n}\sum_{i=1}^n Z_i \xrightarrow{\text{p.s.}} \pi.
\]
\end{exemple}

\subsection{Th\'eor\`eme de Glivenko--Cantelli}

\begin{theoreme}[Glivenko--Cantelli]
Soit $(X_i)$ i.i.d.\ de CDF $F$ et $\hat{F}_n(x) = \frac{1}{n}\sum_{i=1}^n
\mathbf{1}_{X_i \leq x}$ la fonction de r\'epartition empirique. Alors~:
\[
    \sup_{x \in \R} \abs{\hat{F}_n(x) - F(x)} \xrightarrow{\text{p.s.}} 0.
\]
\end{theoreme}

% ------------------------------------------------------------
\section{Convergence en loi (introduction)}
% ------------------------------------------------------------

\begin{definition}[Convergence en loi]
$(X_n)$ \textbf{converge en loi} vers $X$, not\'e $X_n \xrightarrow{\mathcal{L}} X$,
si pour toute fonction $f : \R \to \R$ continue et born\'ee~:
\[
    \lim_{n \to \infty} \E[f(X_n)] = \E[f(X)].
\]
De mani\`ere \'equivalente~: $F_{X_n}(x) \to F_X(x)$ en tout point $x$ de
continuit\'e de $F_X$.
\end{definition}

\begin{theoreme}[Hi\'erarchie compl\`ete]
\[
    X_n \xrightarrow{\text{p.s.}} X \implies X_n \xrightarrow{\PP} X
    \implies X_n \xrightarrow{\mathcal{L}} X.
\]
La convergence en loi est la plus faible. Elle ne concerne que les lois, pas les
variables elles-m\^emes.
\end{theoreme}

\begin{formules}[title=\textbf{R\'esum\'e des convergences}]
\begin{center}
\begin{tabular}{ll}
    \toprule
    \textbf{Convergence} & \textbf{D\'efinition} \\
    \midrule
    Presque s\^ure & $\PP(\lim X_n = X) = 1$ \\
    En probabilit\'e & $\PP(\abs{X_n - X} > \varepsilon) \to 0$ \\
    Dans $L^p$ & $\E[\abs{X_n - X}^p] \to 0$ \\
    En loi & $\E[f(X_n)] \to \E[f(X)]$ pour $f$ continue born\'ee \\
    \bottomrule
\end{tabular}
\end{center}
\end{formules}

% ------------------------------------------------------------
\section{Exercices}
% ------------------------------------------------------------

\begin{exercice}
Soit $X_n \sim \mathcal{U}([0, 1/n])$. Montrer que $X_n \xrightarrow{\PP} 0$
et que $X_n \xrightarrow{L^2} 0$. La convergence est-elle presque s\^ure~?
\end{exercice}

\begin{exercice}
Soit $(X_n)$ i.i.d.\ avec $\PP(X_n = 1) = \PP(X_n = -1) = 1/2$.
Montrer que $\bar{X}_n \to 0$ p.s.\ et que $\PP(\bar{X}_n = 0) \to 0$ quand $n$
est impair. Commenter.
\end{exercice}

\begin{exercice}
Donner un exemple de suite $(X_n)$ convergeant en probabilit\'e vers 0 mais pas
presque s\^urement.
\textit{Indication~:} consid\'erer des indicatrices d'intervalles se d\'epla\c{c}ant
sur $[0,1]$.
\end{exercice}

\begin{exercice}
On tire $n$ nombres ind\'ependamment et uniform\'ement dans $[0,1]$. Soit
$M_n = \max(X_1, \ldots, X_n)$. Montrer que $M_n \to 1$ p.s.
\end{exercice}

\begin{exercice}
Soit $(X_n)$ i.i.d.\ de loi exponentielle $\text{Exp}(1)$. Montrer que
$\frac{1}{n}\max(X_1, \ldots, X_n) \to 0$ p.s.\ et m\^eme que
$\max(X_1,\ldots,X_n)/\ln n \to 1$ p.s.
\end{exercice}

\begin{exercice}
Soit $(X_n)$ i.i.d.\ de loi de Cauchy standard. Montrer que $\bar{X}_n$ ne converge
vers aucune constante en probabilit\'e. \textit{Indication~:} utiliser les fonctions
caract\'eristiques.
\end{exercice}

\begin{exercice}[Monte Carlo]
On veut estimer $I = \int_0^1 e^{-x^2}\, dx$ par Monte Carlo. \'Ecrire
l'estimateur de Monte Carlo et, en utilisant l'in\'egalit\'e de Tchebychev,
d\'eterminer le nombre de tirages n\'ecessaires pour que l'erreur soit inf\'erieure
\`a $0{,}01$ avec probabilit\'e $0{,}95$.
\end{exercice}

\begin{exercice}
Soit $(X_n)$ une suite i.i.d.\ avec $\E[X_1] = 0$ et $\Var(X_1) = 1$. Montrer
par la LFGN que $S_n/n \to 0$ en probabilit\'e, puis que $S_n^2/n^2 \to 0$ en
probabilit\'e. En d\'eduire un r\'esultat sur $S_n/n^{3/4}$.
\end{exercice}
